\section{Introduction}
\label{sec:introduction}

The proliferation of multi-device computing environments has created a
persistent need for seamless content sharing across heterogeneous platforms.
While cloud-based solutions such as AirDrop~\cite{apple2023airdrop},
Google Nearby Share~\cite{google2023nearby}, and Huawei Share
address this need, they typically require vendor lock-in, account
authentication, or explicit pairing protocols. In scenarios
where privacy, latency, and simplicity are paramount---for example, in
enterprise environments behind firewalls, classroom demonstrations, or
offline field deployments---a lightweight, infrastructure-free alternative
is desirable.

Concurrently, advances in hand gesture recognition have made contactless
human--computer interaction increasingly practical. MediaPipe
Hands~\cite{zhang2020mediapipe} provides robust real-time hand landmark
detection on commodity hardware, while lightweight neural architectures
such as Temporal Convolutional Networks (TCNs)~\cite{bai2018empirical,
lea2017temporal} enable accurate temporal pattern recognition with
minimal computational overhead.

In this work, we present \textbf{GrabDrop}, a system that bridges these
two domains by enabling \emph{gesture-driven screenshot sharing} over a
local area network. The key idea is intuitive: a user \emph{grabs} a
screenshot by closing their hand (palm to fist), and another user
\emph{releases} it by opening their hand (fist to palm) on a different
device. The transfer happens instantly over WiFi with no cloud, no
accounts, and no pairing required. Additionally, swipe gestures (hand
moving up or down) enable remote page scrolling.

\subsection{Contributions}

The main contributions of this project are as follows:

\begin{enumerate}[leftmargin=*]
  \item \textbf{End-to-end gesture recognition pipeline.} We design a
    complete pipeline from raw video capture through hand landmark
    extraction (MediaPipe), feature engineering (144~dimensions),
    temporal classification (TCN), and gesture event emission, operating
    in real-time with a two-stage (Idle/Wakeup) power-aware detection
    scheme.

  \item \textbf{Model optimization for edge deployment.} We apply
    structured pruning (30\% channel reduction, yielding
    1.9$\times$~compression from 87K to 46K parameters) followed by
    INT8 static quantization, producing a deployment-ready ONNX model
    of approximately 140\,KB with minimal accuracy degradation.

  \item \textbf{Cross-platform application.} We implement native clients
    for Android (Kotlin/Jetpack~Compose) and Desktop (Python), sharing a
    unified UDP/TCP network protocol for device discovery and screenshot
    transfer. Both platforms support dual detection backends (Neural
    Network and Legacy rule-based) with automatic fallback.

  \item \textbf{Comprehensive data augmentation.} We employ seven
    augmentation techniques including temporal reversal with label
    remapping, expanding a small video dataset into thousands of
    training samples.
\end{enumerate}

The remainder of this report is organized as follows:
Section~\ref{sec:related_work} reviews related work;
Section~\ref{sec:data} describes data collection and preprocessing;
Section~\ref{sec:methodology} details the gesture detection methodology
and TCN architecture;
Section~\ref{sec:training} presents the training procedure;
Section~\ref{sec:deployment} covers model optimization and deployment;
Section~\ref{sec:application} describes the cross-platform application
design;
Section~\ref{sec:results} reports experimental results; and
Section~\ref{sec:conclusion} concludes with future directions.
